\section*{TỰ ĐỘNG HÓA CẤU HÌNH MẠNG DOANH \\NGHIỆP SỬ DỤNG PYTHON \\TÓM TẮT}

\begin{itemize}
    \item[-] \textbf{Vấn đề nghiên cứu:}
    
    Báo cáo gồm 4 chương
        \begin{itemize}
        \item[•] Chương 1 - Tổng quan đề tài
        \item[•] Chương 2 - Cơ sở lý thuyết
        \item[•] Chương 3 - Triển khai hệ thống
        \item[•] Chương 4 - Quy trình tự động hóa
        \item[•] Chương 5 - Kết quả demo
        \end{itemize}
    \item[-] \textbf{Các hướng tiếp cận:}
        \begin{itemize}
        \item[•] Lý thuyết:
        Phân tích sâu các tài liệu khoa học, báo cáo kỹ thuật (Cisco, IETF) về các giao thức mạng cốt lõi (OSPF, RIP, VLAN).
        Nghiên cứu chuyên sâu về Python và thư viện Netmiko để nắm vững cơ chế tương tác và quản lý thiết bị mạng.
        \item[•]Thực hành:
        Xây dựng mô hình mạng ảo hóa trên GNS3-VM sử dụng thiết bị Cisco IOS để mô phỏng môi trường doanh nghiệp.
        Sử dụng Python/Netmiko để code các kịch bản tự động hóa cho từng chức năng (VLAN, OSPF, Quản lý IP).
        Tiến hành kiểm thử từng module, sau đó tích hợp thành một quy trình hoàn chỉnh và đánh giá độ tin cậy so với phương pháp thủ công.
        \end{itemize}
    \item[-] \textbf{Cách giải quyết vấn đề:}
    Xem lại những nội dung đã học qua slide bài giảng, các kiến thức được ghi chép lại trong quá trình học và nghiên cứu thêm các video bài giảng trên mạng. Vận dụng chúng vào để giải quyết các nội dung trong bài báo cáo.
    \item[-] \textbf{Một số kết quả đạt được:}
    Ôn lại được những kiến thức đã học, nắm vững các lý thuyết và phương pháp đã được học trong các môn về mạng máy tính. Rèn luyện tư duy logic cho việc học tập sau này.
\end{itemize}

\newpage